% Part: first-order-logic
% Chapter: model-theory
% Section: partial-iso

\documentclass[../../../include/open-logic-section]{subfiles}

\begin{document}

\olfileid{mod}{bas}{pis}
\section{部分同构}

\begin{defn}
  给定两个结构 $\Struct{M}$ 和 $\Struct{N}$，从 $\Struct{M}$ 到 $\Struct{N}$ 的一个
  \emph{部分同构}是一个有限偏函数 $p$，其自变元取自 $\Domain M$ 且取值于 $\Domain N$，
  并在其定义域上满足 \olref[iso]{defn:isomorphism} 中的同构条件：
  \begin{enumerate}
  \item $p$ 是单射的；
  \item 对每个常项符号~$c$：若 $p(\Assign{c}{M})$ 有定义，
    则 $p(\Assign{c}{M}) = \Assign{c}{N}$；
  \item 对每个 $n$ 元谓词符号 $P$：若 $a_1$、\dots、$a_n$
    在 $p$ 的定义域中，则 $\langle a_1, \dots, a_n\rangle \in
    \Assign P M$ 当且仅当 $\langle p(a_1), \dots, p(a_n) \rangle
    \in \Assign P N$；
  \item 对每个 $n$ 元函数符号 $f$：若 $a_1$、\dots、$a_n$
    在 $p$ 的定义域中，则 $p(\Assign f M (a_1, \dots,a_n))
    = \Assign f N (p(a_1), \dots, p(a_n))$。
  \end{enumerate}
  所谓 $p$ 是有限的，是指 $\dom{p}$ 是有限的。
\end{defn}

注意空函数~$\emptyset$ 总是任意两个结构之间的一个部分同构。

\begin{defn}\ollabel{defn:partialisom}
  两个结构 $\Struct{M}$ 与 $\Struct{N}$ 是\emph{部分同构的}，记作 $\Struct{M} \iso[p]
  \Struct{N}$，
  当且仅当存在一个 $\Struct{M}$ 与 $\Struct{N}$ 之间的非空部分同构集合 $I$，
  满足以下的\emph{来回性质}：
  \begin{enumerate}
  \item (\emph{向前}) 对每个 $p \in I$ 和每个 $a \in \Domain M$，
    存在 $q \in I$ 使得 $p \subseteq q$ 且 $a$ 在 $q$ 的定义域中；
  \item (\emph{向后}) 对每个 $p \in I$ 和每个 $b \in \Domain N$，
    存在 $q \in I$ 使得 $p \subseteq q$ 且 $b$ 在 $q$ 的值域中。
  \end{enumerate}
\end{defn}

\begin{thm}\ollabel{thm:p-isom1}
  如果 $\Struct{M} \iso[p] \Struct{N}$ 且 $\Struct{M}$ 和
  $\Struct{N}$ 是可数的，则 $\Struct{M} \iso
  \Struct{N}$。
\end{thm}

\begin{proof}
  由于 $\Struct{M}$ 和 $\Struct{N}$ 是可数的，设 $\Domain{M} =
  \{a_0, a_1, \ldots \}$ 和 $\Domain{N} = \{b_0, b_1, \ldots \}$。从
  任意一个 $p_0 \in I$ 出发，我们如下定义一个递增的部分同构序列 $p_0 \subseteq p_1 \subseteq p_2
  \subseteq \cdots$：
  \begin{enumerate}
  \item 若 $n+1$ 是奇数，比方说 $n = 2r$，则利用向前性质
    找到一个 $p_{n+1} \in I$，使得 $p_n \subseteq p_{n+1}$
    且 $a_r$ 在 $p_{n+1}$ 的定义域中；
  \item 若 $n+1$ 是偶数，比方说 $n+1 = 2r$，则利用向后性质
    找到一个 $p_{n+1} \in I$，使得 $p_n \subseteq p_{n+1}$
    且 $b_r$ 在 $p_{n+1}$ 的值域中。
  \end{enumerate}
  现令：
  \[
p = \bigcup_{n\ge 0} p_n,
\]
  则 $p$ 是 $\Struct{M}$ 与 $\Struct{N}$ 之间的一个同构。
\end{proof}

\begin{prob}
  详细证明 \olref[mod][bas][pis]{thm:p-isom1} 中定义的 $p$ 确实是一个同构。
\end{prob}

\begin{thm}\ollabel{thm:p-isom2}
  设 $\Struct{M}$ 和 $\Struct{N}$ 是纯关系语言（即只包含
  谓词符号，不含函数符号或常项符号的语言）的结构。那么若
  $\Struct{M} \iso[p] \Struct{N}$，则 $\Struct{M} \elemequiv
  \Struct{N}$。
\end{thm}

\begin{proof}
  通过对公式进行归纳可以证明：若 $a_1$、\dots、$a_n$ 和
  $b_1$、\dots、$b_n$ 使得存在一个部分同构 $p$
  将每个 $a_i$ 映射到 $b_i$，且 $s_1(x_i) =a_i$ 且 $s_2(x_i) =b_i$
  （其中 $i = 1$，\dots，$n$），则 $\Sat{M}{!A}[s_1]$ 当且仅当
  $\Sat{N}{!A}[s_2]$。$n=0$ 时的情形便给出 $\Struct{M} \elemequiv \Struct{N}$。
\end{proof}

\begin{rem}
若语言中包含函数符号，上述结果仍然成立，但需要考察由 $p$ 在
由 $a_1$、\dots、$a_n$ 生成的 $\Struct{M}$ 的子结构与由 $b_1$、\dots、$b_n$ 生成的 $\Struct{N}$ 的子结构之间诱导的同构。
\end{rem}

通过建立公式中嵌套量词的层数与相关部分同构可被扩充的次数之间的联系，上述结果可被“分解”为不同的阶段。

\begin{defn}
  对任意公式 $!A$，其\emph{量词秩}，记作
  $\QuantRank{!A}$（$\QuantRank{!A} \in \Nat$），递归地定义为 $!A$ 中嵌套量词的
  最大层数。两个结构 $\Struct{M}$ 与 $\Struct{N}$ 是\emph{$n$-等价的}，
  记作 $\Struct{M} \elemequiv[n] \Struct{N}$，如果它们在所有量词秩小于等于 $n$ 的语句上真值一致。
\end{defn}

\begin{prop}\ollabel{prop:qr-finite}
  设 $\Lang{L}$ 是一个有限的纯关系语言，即一个包含有限多个谓词符号和常项符号，
  且没有函数符号的语言。则对每个 $n \in \Nat$，在语言
  $\Lang{L}$ 中，量词秩不超过 $n$ 的一阶语句在逻辑等价的意义下只有有限多个。
\end{prop}

\begin{proof}
  对 $n$ 进行归纳。
\end{proof}

\begin{defn}
  给定结构~$\Struct{M}$，令 $\Domain M^{<\omega}$ 为 $\Domain{M}$ 上所有有限序列的集合。我们使用 $\mathbf{a},
  \mathbf{b}, \mathbf{c}, \ldots$ 来表示元素的有限序列。
  如果 $\mathbf{a} \in \Domain{M}^{<\omega}$ 且 $a \in \Domain{M}$，那么 $\mathbf{a}a$ 表示 $\mathbf{a}$ 与 $a$ 的\emph{拼接}。
\end{defn}

\begin{defn}
  给定结构 $\Struct{M}$ 和 $\Struct{N}$，我们按 $n$ 递归定义等长序列之间的关系 $I_n \subseteq \Domain M^{<\omega} \times \Domain N^{<\omega}$ 如下：
   \begin{enumerate}
   \item $I_0(\mathbf{a},\mathbf{b})$ 当且仅当 $\mathbf{a}$ 和 $\mathbf{b}$ 在 $\Struct{M}$ 和 $\Struct{N}$ 中满足相同的原子公式；即，如果 $s_1(x_i) = a_i$ 且 $s_2(x_i) =
     b_i$ 且 $!A$ 是原子公式，且其所有变元都在 $x_1$, \dots,~$x_n$ 之中，那么 $\Sat{M}{!A}[s_1]$ 当且仅当~$\Sat{N}{!A}[s_2]$。
   \item $I_{n+1} (\mathbf{a},\mathbf{b})$ 当且仅当对于每个 $a\in \Domain M$，存在 $b\in \Domain N$ 使得 $I_n
     (\mathbf{a}a,\mathbf{b}b)$，反之亦然。
   \end{enumerate}
\end{defn}


\begin{defn}
  如果 $I_n(\emptyseq,\emptyseq)$ 对 $\Struct{M}$ 和 $\Struct{N}$ 成立（其中 $\emptyseq$ 是空序列），则记作 $\Struct{M} \approx_n \Struct{N}$。
\end{defn}

\begin{thm}\ollabel{thm:b-n-f}
  设 $\Lang{L}$ 是一个纯关系语言。则 $I_n
  (\mathbf{a},\mathbf{b})$ 蕴含：对每个满足 $\QuantRank{!A} \le n$ 的 $!A$，都有 $\Sat{M}{!A}[\mathbf{a}]$ 当且仅当 $\Sat{N}{!A}[\mathbf{b}]$（此处 $\mathbf{a}$ 满足 $!A$ 仍是指任何满足 $s(x_i) = a_i$ 的 $s$ 都满足 $!A$）。此外，若 $\Lang{L}$ 有限，其逆命题也成立。
\end{thm}

\begin{proof}
  $I_n(\mathbf{a},\mathbf{b})$ 蕴含 $\mathbf{a}$ 和 $\mathbf{b}$ 满足相同量词秩不超过 $n$ 的公式，这一点的证明通过对 $!A$ 进行简单的归纳即可完成。对于逆命题，我们通过对 $n$ 的归纳来证明，其中使用了 \olref{prop:qr-finite}，它保证了对于每个 $n$，至多只有有限多个互不等价的公式具有该量词秩。

  对于 $n=0$，假设 $\mathbf{a}$ 和 $\mathbf{b}$ 满足相同的无量词公式，可推出它们满足相同的原子公式，因此有 $I_0(\mathbf{a},\mathbf{b})$。

  对于 $n+1$ 的情况，假设 $\mathbf{a}$ 和 $\mathbf{b}$ 满足相同量词秩不超过 $n+1$ 的公式；为证明 $I_{n+1}(\mathbf{a},\mathbf{b})$，只需证明对于每个 $a \in \Domain M$，存在 $b \in \Domain N$ 使得 $I_n(\mathbf{a}a,\mathbf{b}b)$，而根据归纳假设，只需证明对于每个 $a \in \Domain M$，存在 $b \in \Domain N$ 使得 $\mathbf{a}a$ 和 $\mathbf{b}b$ 满足相同量词秩不超过 $n$ 的公式。

  给定 $a \in \Domain M$，令 $!T^a_n$ 为 $\mathbf{a}a$ 在 $\Struct{M}$ 中满足的、量词秩不超过 $n$ 的所有公式 $!B(x,\mathbf{y})$ 组成的集合；$!T^a_n$ 是有限的，因此我们可以假设它是一个单个的一阶公式。由此可知，$\mathbf{a}$ 满足 $\lexists[x][!T^a_n(x,\mathbf{y})]$，该公式的量词秩不超过 $n+1$。根据假设，$\mathbf{b}$ 在 $\Struct{N}$ 中也满足同一个公式，因此存在 $b \in \Domain N$ 使得 $\mathbf{b}b$ 满足 $!T^a_n$；特别地，$\mathbf{b}b$ 与 $\mathbf{a}a$ 满足相同量词秩不超过 $n$ 的公式。类似地可以证明，对于每个 $b \in \Domain N$，都存在 $a\in \Domain M$ 使得 $\mathbf{a}a$ 和 $\mathbf{b}b$ 满足相同量词秩不超过 $n$ 的公式，这就完成了证明。
\end{proof}

\begin{cor}\ollabel{cor:b-n-f}
  若 $\Struct{M}$ 和 $\Struct{N}$ 是有限语言中的纯关系结构，则 $\Struct{M} \approx_n\Struct{N}$ 当且仅当 $\Struct{M} \elemequiv[n] \Struct{N}$。特别地，$\Struct{M} \elemequiv \Struct{N}$ 当且仅当对于每个 $n$，有 $\Struct{M} \approx_n \Struct{N}$。
\end{cor}

\end{document}